\documentclass[a4paper,12pt]{article}

% -------------------------------------------------
% Кодировка, язык, математика
% -------------------------------------------------
\usepackage[T2A]{fontenc}
\usepackage[utf8]{inputenc}
\usepackage[russian]{babel}
\usepackage{amsmath,amssymb,mathtools}
\usepackage{helvet}
\renewcommand{\familydefault}{\sfdefault}
\usepackage{icomma}

% -------------------------------------------------
% Вёрстка
% -------------------------------------------------
\usepackage[
  a4paper,
  left=20mm,
  right=20mm,
  top=20mm,
  bottom=22mm
]{geometry}

\usepackage{microtype}
\usepackage{parskip}
\usepackage{setspace}
\onehalfspacing
\usepackage{enumitem}
\usepackage{tabularx,booktabs}
\usepackage{xcolor}
\usepackage{tikz}
\usetikzlibrary{calc}
\usepackage{titlesec}
\usepackage{fancyhdr}
\usepackage{hyperref}

% -------------------------------------------------
% Палитра
% -------------------------------------------------
\definecolor{accent}{HTML}{008080}
\definecolor{darkaccent}{HTML}{004D4D}
\definecolor{textgray}{HTML}{333333}
\definecolor{softgray}{HTML}{F2F4F4}
\definecolor{mutedgray}{HTML}{6B7275}

\color{textgray}

% -------------------------------------------------
% Заголовки
% -------------------------------------------------
\titleformat{\section}
  {\Large\bfseries\color{darkaccent}}
  {\thesection.}{0.6em}{}

\titleformat{\subsection}
  {\large\bfseries\color{accent}}
  {\thesubsection.}{0.55em}{}

\titlespacing*{\section}{0pt}{1.5em}{0.7em}
\titlespacing*{\subsection}{0pt}{1.2em}{0.45em}

% -------------------------------------------------
% Колонтитулы
% -------------------------------------------------
\pagestyle{fancy}
\fancyhf{}
\fancyhead[L]{\small\color{mutedgray}Алгебраические преобразования}
\fancyhead[R]{\small\color{mutedgray}Кристина}
\fancyfoot[C]{\small\color{mutedgray}\thepage}
\renewcommand{\headrulewidth}{0.3pt}
\renewcommand{\headrule}{%
  \hbox to\headwidth{\color{accent}\leaders\hrule height \headrulewidth\hfill}}
\setlength{\headheight}{15pt}

% -------------------------------------------------
% Служебные настройки
% -------------------------------------------------
\hypersetup{hidelinks}
\setlength{\parindent}{0pt}
\setlist[itemize]{leftmargin=1.5em,itemsep=0.25em,topsep=0.4em}
\setlist[enumerate]{leftmargin=1.7em,itemsep=0.45em,topsep=0.45em}
\renewcommand{\arraystretch}{1.35}
\raggedbottom

\newcommand{\term}[1]{\textbf{\color{darkaccent}#1}}
\newcommand{\important}[1]{\textbf{\color{accent}#1}}
\newcommand{\separator}{%
  \vspace{0.4em}
  {\color{accent}\hrule height 0.45pt}
  \vspace{0.7em}
}

\begin{document}

% =================================================
% ТИТУЛЬНЫЙ ЛИСТ
% =================================================

\begin{titlepage}
\thispagestyle{empty}

\begin{center}

\vspace*{24mm}

{\small\color{accent}\bfseries УЧЕБНЫЙ ЧЕК-ЛИСТ}

\vspace{10mm}

{\Huge\bfseries\color{darkaccent}
Разложение на множители\\[2mm]
и алгебраические дроби
}

\vspace{8mm}

{\Large
Квадратный трёхчлен, ФСУ,\\
сокращение и преобразование дробей
}

\vspace{18mm}

{\large
Ключевые приёмы для уверенной работы\\
с алгебраическими выражениями
}

\vfill

{\large\bfseries
Лёвин Артём Александрович\\[2mm]
эксклюзивно для Кристины
}

\vspace{10mm}

{\large \today}

\end{center}

% Сдержанная кривая Безье внизу титульного листа
\begin{tikzpicture}[remember picture,overlay]
  \draw[
    accent,
    line width=0.9pt
  ]
  ($(current page.south west)+(18mm,16mm)$)
    .. controls
    ($(current page.south west)+(58mm,27mm)$)
    and
    ($(current page.south east)+(-60mm,7mm)$)
    ..
  ($(current page.south east)+(-18mm,18mm)$);
\end{tikzpicture}

\end{titlepage}

% =================================================
% ОСНОВНОЙ ЧЕК-ЛИСТ
% =================================================

\section{Главная идея занятия}

При работе с алгебраическими дробями полезно постоянно искать возможность
\important{разложить выражения на множители и выполнить сокращение}.

Базовый алгоритм:

\begin{enumerate}
  \item Найдите квадратные трёхчлены.
  \item Проверьте возможность вынесения общего множителя.
  \item Проверьте формулы сокращённого умножения.
  \item Разложите выражения на множители.
  \item Сравните получившиеся множители в числителе и знаменателе.
  \item При необходимости преобразуйте знаки с помощью вынесения минуса.
  \item Выполните допустимые сокращения.
  \item Сохраните все ограничения исходного выражения.
\end{enumerate}

\separator

\term{Главный ориентир:}
после каждого преобразования задавайте себе вопрос:
\emph{«Появилась ли возможность сократить или упростить выражение?»}

% -------------------------------------------------

\section{Квадратный трёхчлен}

Квадратный трёхчлен имеет вид

\[
ax^2+bx+c,\qquad a\ne0.
\]

Если его корни равны \(x_1\) и \(x_2\), то

\[
\boxed{
ax^2+bx+c=a(x-x_1)(x-x_2)
}
\]

Коэффициент \(a\) является частью формулы и сохраняется при разложении.
При \(a=1\) множитель \(1\) обычно отдельно не записывают.

\subsection{Поиск корней через дискриминант}

\[
D=b^2-4ac.
\]

При \(D>0\):

\[
x_{1,2}=\frac{-b\pm\sqrt D}{2a}.
\]

При \(D=0\) существует один различный корень кратности \(2\):

\[
x_1=x_2=-\frac{b}{2a}.
\]

Тогда

\[
ax^2+bx+c=a(x-x_1)^2.
\]

При \(D<0\) действительных корней нет, поэтому над множеством действительных
чисел разложение на линейные множители такого вида отсутствует.

\subsection{Пример 1}

Разложим на множители

\[
x^2-11x+28.
\]

Здесь

\[
a=1,\qquad b=-11,\qquad c=28.
\]

\[
D=(-11)^2-4\cdot1\cdot28
=121-112=9.
\]

\[
\sqrt D=3.
\]

\[
x_1=\frac{-(-11)+3}{2}
=\frac{14}{2}=7,
\]

\[
x_2=\frac{-(-11)-3}{2}
=\frac{8}{2}=4.
\]

Следовательно,

\[
\boxed{
x^2-11x+28=(x-7)(x-4)
}
\]

\term{Контроль знаков:}
если \(b<0\), в формуле корней возникает конструкция
\(-b=-(-11)\). Двойной минус требуется записывать явно.

\subsection{Пример 2}

\[
x^2-x-6.
\]

\[
D=(-1)^2-4\cdot1\cdot(-6)
=1+24=25.
\]

\[
x_1=3,\qquad x_2=-2.
\]

Поэтому

\[
\boxed{
x^2-x-6=(x-3)(x+2)
}
\]

Знак \(+\) во второй скобке возникает из выражения

\[
x-(-2)=x+2.
\]

\subsection{Отрицательный старший коэффициент}

Рассмотрим

\[
-x^2+3x+4.
\]

Здесь

\[
a=-1,\qquad b=3,\qquad c=4.
\]

\[
D=3^2-4\cdot(-1)\cdot4=25.
\]

\[
x_1=-1,\qquad x_2=4.
\]

Следовательно,

\[
-x^2+3x+4
=-(x+1)(x-4).
\]

Внешний минус можно внести в одну из скобок:

\[
-(x-4)=4-x.
\]

Поэтому удобная форма:

\[
\boxed{
-x^2+3x+4=(x+1)(4-x)
}
\]

Именно такая запись особенно полезна при сокращении дробей.

% -------------------------------------------------

\section{Формулы сокращённого умножения}

Три наиболее важные формулы для текущих задач:

\[
\boxed{
A^2-B^2=(A-B)(A+B)
}
\]

\[
\boxed{
A^2-2AB+B^2=(A-B)^2
}
\]

\[
\boxed{
A^2+2AB+B^2=(A+B)^2
}
\]

\subsection{Примеры}

Разность квадратов:

\[
9y^2-16
=(3y)^2-4^2
=\boxed{(3y-4)(3y+4)}.
\]

Квадрат разности:

\[
x^2-4x+4
=x^2-2\cdot x\cdot2+2^2
=\boxed{(x-2)^2}.
\]

Более общий вид:

\[
9x^2-6xy+y^2
=(3x)^2-2\cdot3x\cdot y+y^2
=\boxed{(3x-y)^2}.
\]

\term{Самопроверка:}
после сворачивания выражения полезно мысленно раскрыть получившийся квадрат
обратно и проверить все три слагаемых.

% -------------------------------------------------

\section{Сокращение алгебраических дробей}

Сокращать можно \important{общие множители} числителя и знаменателя.

Например,

\[
\frac{(x-2)(x+1)}
     {(x-2)(x+4)}
=
\frac{x+1}{x+4},
\]

при этом из исходного выражения сохраняются ограничения

\[
x\ne2,\qquad x\ne-4.
\]

После сокращения ограничение \(x\ne2\) продолжает действовать.

\subsection{Слагаемые сокращать нельзя}

Выражение

\[
\frac{x+2}{x+5}
\]

сократить нельзя.

Также сокращение невозможно в конструкции

\[
\frac{x+2}{x(x+2)+1},
\]

поскольку \(x+2\) здесь не является множителем всего знаменателя.

\term{Практическое правило:}
перед сокращением выражение должно быть представлено как произведение
множителей.

% -------------------------------------------------

\section{Вынесение и внесение минуса}

Противоположные выражения связаны равенством

\[
A-B=-(B-A).
\]

Например,

\[
4-x=-(x-4).
\]

Это позволяет преобразовывать знаменатели и получать одинаковые множители.

Для дроби справедливо

\[
-\frac{A}{B}
=
\frac{-A}{B}
=
\frac{A}{-B},
\qquad B\ne0.
\]

Также

\[
\frac{-A}{-B}=\frac{A}{B}.
\]

Скобки здесь обязательны, когда минус относится ко всему выражению:

\[
-(3x-2)=-3x+2.
\]

\subsection{Пример сокращения}

Упростим

\[
\frac{-x^2+3x+4}{4-x}.
\]

Числитель уже раскладывался:

\[
-x^2+3x+4=(x+1)(4-x).
\]

Поэтому

\[
\frac{(x+1)(4-x)}{4-x}
=x+1,
\]

при условии

\[
x\ne4.
\]

Итак,

\[
\boxed{x+1},\qquad x\ne4.
\]

% -------------------------------------------------

\section{Общий знаменатель}

Перед приведением дробей к общему знаменателю полезно сначала разложить
знаменатели на множители.

Пример:

\[
\frac{1}{x-2}+\frac{2}{x+2}.
\]

Общий знаменатель:

\[
(x-2)(x+2).
\]

Поэтому

\[
\frac{1}{x-2}+\frac{2}{x+2}
=
\frac{x+2+2(x-2)}
     {(x-2)(x+2)}.
\]

Раскрываем скобки:

\[
x+2+2x-4=3x-2.
\]

Получаем

\[
\boxed{
\frac{3x-2}{x^2-4}
},
\qquad
x\ne\pm2.
\]

\term{Важно:}
домножаемое выражение записывайте в скобках. Это особенно существенно
при вычитании дробей.

Например,

\[
A-(B-C)=A-B+C.
\]

Скобки защищают знаки от потери при раскрытии.

% -------------------------------------------------

\section{Деление алгебраических дробей}

При делении вторую дробь переворачиваем и заменяем деление умножением:

\[
\frac{A}{B}:\frac{C}{D}
=
\frac{A}{B}\cdot\frac{D}{C}.
\]

Необходимо учитывать

\[
B\ne0,\qquad D\ne0,\qquad C\ne0.
\]

Последнее условие возникает потому, что делитель не может быть равен нулю.

\subsection{Пример}

\[
\frac{3x}{x-2}:\frac{x}{x+1}.
\]

Переходим к умножению:

\[
\frac{3x}{x-2}\cdot\frac{x+1}{x}.
\]

Сокращаем множитель \(x\):

\[
\boxed{
\frac{3(x+1)}{x-2}
}.
\]

Ограничения исходного выражения:

\[
\boxed{
x\ne2,\qquad x\ne-1,\qquad x\ne0.
}
\]

% -------------------------------------------------

\section{Когда алгебраическая дробь равна нулю}

Для уравнения

\[
\frac{P(x)}{Q(x)}=0
\]

используем систему

\[
\boxed{
\begin{cases}
P(x)=0,\\
Q(x)\ne0.
\end{cases}
}
\]

Сначала находим корни числителя, затем исключаем значения,
обращающие знаменатель в нуль.

\subsection{Пример}

Решим

\[
\frac{x^2-x-6}{x+2}=0.
\]

Разложим числитель:

\[
x^2-x-6=(x-3)(x+2).
\]

Получаем систему

\[
\begin{cases}
(x-3)(x+2)=0,\\
x+2\ne0.
\end{cases}
\]

Из первого условия:

\[
x=3
\quad\text{или}\quad
x=-2.
\]

Из второго:

\[
x\ne-2.
\]

Следовательно,

\[
\boxed{x=3}.
\]

% -------------------------------------------------

\section{Типичные ошибки}

\begin{enumerate}
  \item Потеря коэффициента \(a\) при разложении квадратного трёхчлена.
  \item Ошибка в знаке коэффициента \(b\).
  \item Потеря минуса в конструкции \(-b\).
  \item Запись \(\sqrt D\) без предварительного вычисления самого \(D\).
  \item Попытка сократить отдельные слагаемые.
  \item Потеря ограничений после сокращения общего множителя.
  \item Неверное внесение внешнего минуса сразу в несколько множителей.
  \item Отсутствие скобок при вычитании сложного выражения.
  \item Потеря условия, что делитель при делении дробей должен быть ненулевым.
  \item Подстановка запрещённого значения в ответ дробного уравнения.
  \item Ранний переход к десятичным приближениям вместо точных обыкновенных дробей.
\end{enumerate}

% -------------------------------------------------

\section{Короткий контрольный алгоритм}

Получив новое алгебраическое выражение, последовательно проверьте:

\begin{enumerate}
  \item Есть ли общий множитель?
  \item Есть ли квадратный трёхчлен?
  \item Видна ли формула сокращённого умножения?
  \item Можно ли преобразовать противоположные скобки вынесением минуса?
  \item Какие ограничения задают знаменатели?
  \item После разложения появились ли одинаковые множители?
  \item Можно ли сократить?
  \item Осталось ли выражение, которое можно дополнительно упростить?
\end{enumerate}

\separator

\begin{center}
\large\bfseries\color{darkaccent}
Разложить на множители
\(\longrightarrow\)
проверить знаки
\(\longrightarrow\)
сократить
\(\longrightarrow\)
проверить ограничения
\end{center}

% =================================================
% ТРЕНИРОВОЧНЫЙ БЛОК
% =================================================

\clearpage

\section*{Тренировочный блок}

\textcolor{mutedgray}{
Выполняйте преобразования подробно.
Все ограничения исходных выражений сохраняйте до конца решения.
}

\separator

\begin{enumerate}[label=\textbf{\arabic*.}]

\item
Разложите на множители квадратный трёхчлен:

\[
2x^2-x-6.
\]

\vspace{12mm}

\item
Упростите дробь и укажите допустимые значения переменной:

\[
\frac{-x^2+5x-6}{3-x}.
\]

\vspace{15mm}

\item
Упростите выражение и укажите все ограничения:

\[
\frac{9x^2-12x+4}{3x-2}
\cdot
\frac{x+1}{6x^2+2x-4}.
\]

\vspace{18mm}

\item
Приведите к общему знаменателю и упростите:

\[
\frac{2}{x-1}
+
\frac{3}{x+1}.
\]

Укажите ограничения.

\vspace{18mm}

\item
Решите уравнение:

\[
\frac{2x^2-7x+3}{x^2-x-6}=0.
\]

Обязательно запишите ограничения, возникающие из знаменателя.

\end{enumerate}

% =================================================
% ОТВЕТЫ
% =================================================

\clearpage

\section*{Ответы к тренировочному блоку}

\renewcommand{\arraystretch}{1.5}

\begin{table}[ht]
\centering
\caption{Ответы}
\begin{tabularx}{\linewidth}{
  >{\centering\arraybackslash}p{12mm}
  X
  >{\raggedright\arraybackslash}p{52mm}
}
\toprule
\textbf{№} &
\textbf{Ответ} &
\textbf{Ограничения / замечание}
\\
\midrule

1 &
\[
2x^2-x-6=(x-2)(2x+3)
\]
&
Проверка раскрытием скобок.
\\

2 &
\[
x-2
\]
&
\[
x\ne3
\]
\\

3 &
\[
\frac12
\]
&
\[
x\ne\frac23,\qquad x\ne-1
\]
\\

4 &
\[
\frac{5x-1}{x^2-1}
\]
&
\[
x\ne-1,\qquad x\ne1
\]
\\

5 &
\[
x=\frac12
\]
&
\[
x\ne3,\qquad x\ne-2
\]
Корень \(x=3\) числителя запрещён знаменателем.
\\

\bottomrule
\end{tabularx}
\end{table}

\vfill

\begin{center}
{\color{accent}\rule{0.32\linewidth}{0.5pt}}

\vspace{4mm}

\textbf{\color{darkaccent}Финальная проверка}

Разложили на множители. Проверили знаки.\\
Выполнили только допустимые сокращения.\\
Сохранили ограничения исходного выражения.
\end{center}

\end{document}